% TOP_PROBLEM: {"id": 1235, "title": "The Sunflower Conjecture", "category_id": 2, "category": {"id": 2, "name": "combinatorics", "display_name": "Combinatorics", "description": "Counting problems, graph theory, discrete structures.", "slug": "combinatorics", "order_index": 2, "created_at": "2026-07-31T15:26:25.671Z"}, "status": "open"}
% FIRST_POSTED: null
% ENTRY_KIND: statement_only
% Imported from: batch05.tex; UnsolvedMath ID 1235
\documentclass[11pt]{article}
\usepackage[margin=25mm]{geometry}
\usepackage{fontspec}
\setmainfont{Latin Modern Roman}
\usepackage{amsmath,amssymb,mathtools}
\usepackage{unicode-math}
\setmathfont{Latin Modern Math}
\usepackage{xurl,hyperref}
\hypersetup{colorlinks=true,urlcolor=blue}
\setcounter{secnumdepth}{0}
\setlength{\parindent}{0pt}
\setlength{\parskip}{5pt}
\setlength{\emergencystretch}{3em}
\newcommand{\sref}[2]{\hyperlink{source:#1:#2}{[S#2]}}
\newcommand{\cataloguescope}[1]{\paragraph{Recorded target.}#1}
\begin{document}
% UnsolvedMath ID 1235; source catalogue code COMB-012
\setcounter{section}{1234}
\section{The Sunflower Conjecture}\label{1235}
{\small\sffamily UnsolvedMath ID 1235 · Combinatorics}
\subsection{Definitions and mathematical statement}

\paragraph{Shared notation.}$\mathbb N=\{1,2,\ldots\}$, and $\mathbb Z,\mathbb Q,\mathbb R,\mathbb C$ denote the integers, rationals, reals and complex numbers. Any other conventions are specified locally.


Let $k\geq1$ and let $\mathcal F$ be a finite family of distinct $k$-element subsets of an arbitrary finite set. Three distinct members $A_1,A_2,A_3$ form a three-petal sunflower if
\[
A_1\cap A_2=A_1\cap A_3=A_2\cap A_3.
\]
The common intersection is the core; the sets obtained by removing it are pairwise disjoint.
\paragraph{Statement recorded in the supplied source.}
Does there exist an absolute constant $c>0$ such that, for every $k\geq1$ and every such family $\mathcal F$, the inequality
\[
|\mathcal F|\geq c^k k!
\]
forces a three-petal sunflower? \sref{1235}{2}
The factorial-bearing assertion above has an affirmative classical answer: the Erd\H{o}s--Rado sunflower lemma gives the bound $k!2^k$ for three petals, so $c=3$ certainly suffices. The stronger exponential sunflower conjecture removes the factor $k!$; that different assertion is not the recorded target here. \sref{1235}{2}
\subsection{Short English statement}
Does a factorial-times-exponential number of distinct sets of the same size force three sets with identical pairwise intersections?
\subsection{Sources}
\begin{description}
\item[\hypertarget{source:1235:1}{S1}] ULAM.AI and the UnsolvedMath Contributors, UnsolvedMath: A Curated Collection of Open Mathematics Problems, record 1235 (COMB-012). CC BY 4.0. \url{https://huggingface.co/datasets/ulamai/UnsolvedMath}.
\item[\hypertarget{source:1235:2}{S2}] Ryan Alweiss, Shachar Lovett, Kewen Wu and Jiapeng Zhang, \emph{Improved bounds for the sunflower lemma}, arXiv:1908.08483v3 (2021), Definition 1.1, Lemma 1.2 and Conjecture 1.3. \url{https://arxiv.org/pdf/1908.08483}.
\end{description}
\label{end:1235}
\end{document}
